%% ma: L12-B16-0008
%% lop: 12
%% chuong: 5
%% bai: 16
%% dang: 12.16.6
%% loai: DS
%% muc: ['TH', 'TH', 'VD', 'VD']
%% dap_an: "SĐĐĐ"
%% nguon: "(NBV)-10 ĐỀ KIỂM TRA HK2-TOÁN 12-DUNG CHUNG 3 SÁCH.pdf (H03, Đề số 3), Phần 2 Câu 1"
%% kien_thuc: "Bài 15 (phương trình đường thẳng), Bài 14 (vectơ pháp tuyến), Bài 16 (góc giữa hai đường thẳng, góc giữa đường thẳng và trục), vị trí tương đối đường thẳng và mặt phẳng"
%% trang_thai: cho-duyet
%% ly_do: "Dạng toán mới đề xuất 12.16.6, chờ giáo viên duyệt"
%% ghi_chu: "Đáp án gốc khớp (SĐĐĐ). Lời giải gốc ghi điểm M(1;3;2) do mất dấu trừ; đúng là M(-1;3;2)."
\begin{ex}
Trong không gian $Oxyz$, cho hai đường thẳng
\[d_1\colon\dfrac{x+2}{-1}=\dfrac{y-1}{-2}=\dfrac{z+3}{2},\qquad d_2\colon\begin{cases}x=-1+t\\y=3+t\\z=2-mt\end{cases}\]
và mặt phẳng $(P)\colon2x+2y+z-3=0$.
\choiceTF
{Khi $m=0$, số đo góc giữa hai đường thẳng $d_1$ và $d_2$ bằng $60^\circ$.}
{\True $\cos\left(d_1,Ox\right)=\dfrac13$.}
{\True Đường thẳng $\Delta$ đi qua gốc tọa độ $O$ và vuông góc với $(P)$ tạo với đường thẳng $d_1$ một góc $\alpha$ có $\cos\alpha=\dfrac49$.}
{\True $d_2$ song song với $(P)$ khi và chỉ khi $m=4$.}
\loigiai{$d_1$ có vectơ chỉ phương $\vec u_1=(-1;-2;2)$, $|\vec u_1|=3$.

a) Khi $m=0$, $d_2$ có $\vec u_2=(1;1;0)$: $\cos(d_1,d_2)=\dfrac{|-1-2|}{3\sqrt2}=\dfrac{\sqrt2}{2}$, nên góc bằng $45^\circ$. Mệnh đề sai.

b) $Ox$ có vectơ chỉ phương $\vec i=(1;0;0)$: $\cos(d_1,Ox)=\dfrac{|-1|}{3\cdot1}=\dfrac13$. Đúng.

c) $\Delta\perp(P)$ nên $\Delta$ có vectơ chỉ phương $\vec n=(2;2;1)$, $|\vec n|=3$: $\cos\alpha=\dfrac{|-2-4+2|}{3\cdot3}=\dfrac49$. Đúng.

d) $d_2$ có $\vec u_2=(1;1;-m)$. $d_2\parallel(P)$ khi $\vec u_2\cdot\vec n=0\Leftrightarrow2+2-m=0\Leftrightarrow m=4$ và điểm $M(-1;3;2)\in d_2$ không thuộc $(P)$ (vì $-2+6+2-3=3\neq0$). Vậy $d_2\parallel(P)\Leftrightarrow m=4$. Đúng.}
\end{ex}
